\chapter{Mathematical Preliminaries}
\label{ch:preliminaries}

This chapter establishes the mathematical foundations required for rigorous adaptive control analysis. We present complete proofs of key results, including Barbalat's Lemma, develop the theory of persistent excitation, and establish the function space framework for systems with memory.

\section{Function Spaces and Norms}

\subsection{Lebesgue Spaces}

\begin{definition}[$\Lp{p}$ Spaces]
For $p \in [1, \infty)$, the Lebesgue space $\Lp{p}[0,\infty)$ consists of all measurable functions $f: [0,\infty) \to \R^n$ satisfying
\begin{equation}
    \norm{f}_{\Lp{p}} = \left( \int_0^\infty \norm{f(t)}^p \, dt \right)^{1/p} < \infty.
\end{equation}
For $p = \infty$, we define
\begin{equation}
    \norm{f}_{\Lp{\infty}} = \esssup_{t \geq 0} \norm{f(t)} < \infty.
\end{equation}
\end{definition}

\begin{remark}
In adaptive control, $\Lp{2}$ is the energy space (finite energy signals), while $\Lp{\infty}$ captures bounded signals. A key question is whether $\Lp{2}$ implies $\Lp{\infty}$ for system states.
\end{remark}

\subsection{Sobolev Spaces}

\begin{definition}[Sobolev Space $W^{k,p}$]
For $k \in \N$ and $p \in [1,\infty]$, the Sobolev space $W^{k,p}[0,\infty)$ consists of functions $f \in \Lp{p}$ whose weak derivatives up to order $k$ also belong to $\Lp{p}$:
\begin{equation}
    W^{k,p} = \left\{ f \in \Lp{p} : D^\alpha f \in \Lp{p} \text{ for all } |\alpha| \leq k \right\},
\end{equation}
with norm
\begin{equation}
    \norm{f}_{W^{k,p}} = \sum_{|\alpha| \leq k} \norm{D^\alpha f}_{\Lp{p}}.
\end{equation}
\end{definition}

\begin{remark}
For adaptive systems, we often require $\dot{x} \in \Lp{2}$ (in addition to $x \in \Lp{2}$), which places $x$ in $W^{1,2}$, the first-order Sobolev space. This regularity is crucial for applying Barbalat's Lemma.
\end{remark}

\subsection{Uniformly Continuous Functions}

\begin{definition}[Uniform Continuity]
A function $f: [0,\infty) \to \R^n$ is \textbf{uniformly continuous} if for every $\epsilon > 0$, there exists $\delta > 0$ such that
\begin{equation}
    \norm{t_1 - t_2} < \delta \implies \norm{f(t_1) - f(t_2)} < \epsilon
\end{equation}
for all $t_1, t_2 \in [0, \infty)$.
\end{definition}

\begin{lemma}[Sufficient Condition for Uniform Continuity]
\label{lem:unif_cont}
If $f: [0,\infty) \to \R^n$ is differentiable and $\dot{f}$ is bounded (i.e., $f \in W^{1,\infty}$), then $f$ is uniformly continuous.
\end{lemma}

\begin{proof}
Assume $\norm{\dot{f}(t)} \leq M$ for all $t \geq 0$. For any $t_1, t_2 \geq 0$, the fundamental theorem of calculus gives
\begin{equation}
    f(t_2) - f(t_1) = \int_{t_1}^{t_2} \dot{f}(s) \, ds.
\end{equation}
Taking norms:
\begin{equation}
    \norm{f(t_2) - f(t_1)} \leq \int_{t_1}^{t_2} \norm{\dot{f}(s)} \, ds \leq M |t_2 - t_1|.
\end{equation}
Given $\epsilon > 0$, choose $\delta = \epsilon / M$. Then $|t_2 - t_1| < \delta$ implies
\begin{equation}
    \norm{f(t_2) - f(t_1)} < M \cdot \frac{\epsilon}{M} = \epsilon,
\end{equation}
establishing uniform continuity.
\end{proof}

\section{Barbalat's Lemma}

Barbalat's Lemma is the cornerstone of adaptive control convergence analysis. It provides conditions under which an $\Lp{2}$ signal must tend to zero.

\begin{lemma}[Barbalat's Lemma]
\label{lem:barbalat}
Let $f: [0,\infty) \to \R$ be a uniformly continuous function. If $f \in \Lp{1}[0,\infty)$ (i.e., $\int_0^\infty |f(t)| \, dt < \infty$), then
\begin{equation}
    \lim_{t \to \infty} f(t) = 0.
\end{equation}
\end{lemma}

\begin{proof}
We prove by contradiction. Suppose $f$ does not converge to zero. Then there exists $\epsilon_0 > 0$ and a sequence $\{t_k\}_{k=1}^\infty$ with $t_k \to \infty$ such that
\begin{equation}
    |f(t_k)| \geq \epsilon_0 \quad \text{for all } k.
\end{equation}

By uniform continuity, there exists $\delta > 0$ such that
\begin{equation}
    |t - t_k| < \delta \implies |f(t) - f(t_k)| < \frac{\epsilon_0}{2}.
\end{equation}

Consider the interval $I_k = [t_k - \delta/2, t_k + \delta/2]$. For $t \in I_k$:
\begin{equation}
    |f(t)| \geq |f(t_k)| - |f(t) - f(t_k)| \geq \epsilon_0 - \frac{\epsilon_0}{2} = \frac{\epsilon_0}{2}.
\end{equation}

Thus, on each interval $I_k$ (which has length $\delta$), we have $|f(t)| \geq \epsilon_0/2$. This gives
\begin{equation}
    \int_{I_k} |f(t)| \, dt \geq \frac{\epsilon_0}{2} \cdot \delta = \frac{\epsilon_0 \delta}{2}.
\end{equation}

Since $t_k \to \infty$, we can choose the sequence such that the intervals $I_k$ are disjoint for sufficiently large $k$. Summing over infinitely many disjoint intervals:
\begin{equation}
    \int_0^\infty |f(t)| \, dt \geq \sum_{k=K}^\infty \int_{I_k} |f(t)| \, dt \geq \sum_{k=K}^\infty \frac{\epsilon_0 \delta}{2} = \infty,
\end{equation}
contradicting $f \in \Lp{1}$. Therefore, $f(t) \to 0$ as $t \to \infty$.
\end{proof}

\subsection{Barbalat's Lemma: Key Corollary}

The following corollary is the form most commonly used in adaptive control.

\begin{corollary}[Barbalat for $\Lp{2}$ Signals]
\label{cor:barbalat_L2}
Let $f: [0,\infty) \to \R^n$ satisfy:
\begin{enumerate}
    \item $f \in \Lp{2}[0,\infty)$ (i.e., $\int_0^\infty \norm{f(t)}^2 \, dt < \infty$),
    \item $\dot{f} \in \Lp{\infty}[0,\infty)$ (i.e., $\dot{f}$ is bounded).
\end{enumerate}
Then $\lim_{t \to \infty} f(t) = 0$.
\end{corollary}

\begin{proof}
\textbf{Step 1: Uniform continuity.}
By Lemma \ref{lem:unif_cont}, condition (2) implies $f$ is uniformly continuous.

\textbf{Step 2: $\Lp{2}$ implies $\Lp{1}$ for uniformly continuous functions.}
We show that $\norm{f} \in \Lp{1}$. Define $g(t) = \norm{f(t)}$. Since $f \in \Lp{2}$:
\begin{equation}
    \int_0^\infty g(t)^2 \, dt < \infty.
\end{equation}

We claim $\int_0^\infty g(t) \, dt < \infty$. To see this, decompose the integral:
\begin{align}
    \int_0^\infty g(t) \, dt &= \int_0^T g(t) \, dt + \int_T^\infty g(t) \, dt.
\end{align}

For any finite $T$, the first integral is finite. For the second integral, we use Cauchy-Schwarz:
\begin{align}
    \int_T^\infty g(t) \, dt &= \int_T^\infty g(t) \cdot 1 \, dt \\
    &\leq \left( \int_T^\infty g(t)^2 \, dt \right)^{1/2} \left( \int_T^\infty 1^2 \, dt \right)^{1/2}.
\end{align}

However, this approach fails since $\int_T^\infty 1 \, dt = \infty$. Instead, we use a different argument.

\textbf{Step 2 (Revised): Direct application to $g^2$.}
Actually, we apply Barbalat directly to $g^2(t) = \norm{f(t)}^2$. We have $g^2 \in \Lp{1}$ by assumption (1). We need to show $g^2$ is uniformly continuous.

Since $f$ is uniformly continuous (each component), for any $\epsilon > 0$, there exists $\delta > 0$ such that
\begin{equation}
    |t_1 - t_2| < \delta \implies \norm{f(t_1) - f(t_2)} < \sqrt{\epsilon/(2M)},
\end{equation}
where $M = \sup_{t \geq 0} \norm{f(t)}$ (which exists since $f \in \Lp{2}$ and $\dot{f}$ bounded implies $f$ is bounded).

Then:
\begin{align}
    |g^2(t_1) - g^2(t_2)| &= |\norm{f(t_1)}^2 - \norm{f(t_2)}^2| \\
    &= |(\norm{f(t_1)} - \norm{f(t_2)})(\norm{f(t_1)} + \norm{f(t_2)})| \\
    &\leq \norm{f(t_1) - f(t_2)} \cdot (\norm{f(t_1)} + \norm{f(t_2)}) \\
    &< \sqrt{\frac{\epsilon}{2M}} \cdot 2M = \sqrt{2M\epsilon}.
\end{align}

Choosing $\epsilon$ small enough makes $g^2$ uniformly continuous. By Barbalat's Lemma, $g^2(t) \to 0$, hence $g(t) = \norm{f(t)} \to 0$.
\end{proof}

\begin{remark}[Practical Implication]
In adaptive control, we often show:
\begin{itemize}
    \item $V(t)$ (Lyapunov function) is bounded $\implies$ states $x, \thetatilde \in \Lp{\infty}$.
    \item $\int_0^\infty V(t) \, dt < \infty$ $\implies$ tracking error $e \in \Lp{2}$.
    \item Closed-loop dynamics $\implies$ $\dot{e}$ is bounded (depends on bounded states).
\end{itemize}
Then Corollary \ref{cor:barbalat_L2} gives $e(t) \to 0$.
\end{remark}

\section{Persistent Excitation}

Persistent excitation (PE) is the key condition distinguishing parameter convergence from mere tracking convergence.

\subsection{Definition and Motivation}

\begin{definition}[Persistent Excitation]
\label{def:pe}
A function $\phi: [0,\infty) \to \R^{n \times m}$ is \textbf{persistently exciting} (PE) if there exist constants $T > 0$, $\alpha > 0$ such that
\begin{equation}
    \int_t^{t+T} \phi(\tau)^\top \phi(\tau) \, d\tau \geq \alpha I_m \quad \text{for all } t \geq 0,
\end{equation}
where the inequality is in the sense of positive definiteness.
\end{definition}

\begin{remark}[Intuition]
PE requires that $\phi$ contains sufficiently rich spectral content over every time window $[t, t+T]$. The matrix $\int_t^{t+T} \phi(\tau)^\top \phi(\tau) \, d\tau$ must be "sufficiently positive definite" uniformly in $t$.
\end{remark}

\begin{example}[Sinusoidal Signal]
Let $\phi(t) = [\sin(\omega t), \cos(\omega t)]^\top$ for $\omega > 0$. Then:
\begin{align}
    \int_t^{t+T} \phi(\tau)^\top \phi(\tau) \, d\tau &= \int_t^{t+T} \begin{bmatrix} \sin^2(\omega \tau) & \sin(\omega \tau)\cos(\omega \tau) \\ \sin(\omega \tau)\cos(\omega \tau) & \cos^2(\omega \tau) \end{bmatrix} d\tau.
\end{align}

Using $\int_t^{t+T} \sin^2(\omega \tau) \, d\tau = T/2$ (for $T = 2\pi/\omega$) and $\int_t^{t+T} \sin(\omega \tau)\cos(\omega \tau) \, d\tau = 0$, we get:
\begin{equation}
    \int_t^{t+T} \phi(\tau)^\top \phi(\tau) \, d\tau = \frac{T}{2} I_2.
\end{equation}

Thus, $\phi$ is PE with $\alpha = T/2 = \pi/\omega$.
\end{example}

\begin{example}[Non-PE Signal]
Consider $\phi(t) = e^{-t} \sin(t)$. As $t \to \infty$, $\phi(t) \to 0$, so
\begin{equation}
    \int_t^{t+T} \phi(\tau)^2 \, d\tau \to 0 \quad \text{as } t \to \infty.
\end{equation}
This violates PE since $\alpha$ cannot be chosen uniformly for all $t$.
\end{example}

\subsection{Properties of PE Signals}

\begin{lemma}[PE and Boundedness]
If $\phi: [0,\infty) \to \R^{n \times m}$ is PE with constants $(T, \alpha)$ and $\phi \in \Lp{\infty}$, then:
\begin{equation}
    \alpha \leq T \cdot \norm{\phi}_{\Lp{\infty}}^2.
\end{equation}
\end{lemma}

\begin{proof}
From the PE condition:
\begin{equation}
    \alpha I_m \leq \int_t^{t+T} \phi(\tau)^\top \phi(\tau) \, d\tau \leq \int_t^{t+T} \norm{\phi(\tau)}^2 I_m \, d\tau \leq T \norm{\phi}_{\Lp{\infty}}^2 I_m.
\end{equation}
Taking the minimum eigenvalue of both sides yields the result.
\end{proof}

\begin{theorem}[PE Implies Parameter Convergence]
\label{thm:pe_convergence}
Consider the parameter error dynamics
\begin{equation}
    \dot{\thetatilde}(t) = -\Gamma \phi(t)^\top e(t),
\end{equation}
where $\Gamma \succ 0$, $\phi$ is PE, and $e \in \Lp{2} \cap \Lp{\infty}$ with $\dot{e} \in \Lp{\infty}$.

Then $\thetatilde(t) \to 0$ as $t \to \infty$ (exponentially).
\end{theorem}

\begin{proof}
We provide a detailed proof in several steps.

\textbf{Step 1: Lyapunov function.}
Define $V(\thetatilde) = \frac{1}{2} \thetatilde^\top \Gamma^{-1} \thetatilde$. Then:
\begin{equation}
    \dot{V} = \thetatilde^\top \Gamma^{-1} \dot{\thetatilde} = -\thetatilde^\top \phi^\top e = -e^\top \phi \thetatilde.
\end{equation}

\textbf{Step 2: Integral inequality.}
Integrate over $[t, t+T]$:
\begin{align}
    V(t) - V(t+T) &= -\int_t^{t+T} e(\tau)^\top \phi(\tau) \thetatilde(\tau) \, d\tau \\
    &\geq -\int_t^{t+T} \norm{e(\tau)} \norm{\phi(\tau)} \norm{\thetatilde(\tau)} \, d\tau.
\end{align}

Since $\thetatilde \in \Lp{\infty}$ (from $V$ bounded), let $M = \sup_{t \geq 0} \norm{\thetatilde(t)}$. Then:
\begin{equation}
    V(t) - V(t+T) \geq -M \int_t^{t+T} \norm{e(\tau)} \norm{\phi(\tau)} \, d\tau.
\end{equation}

\textbf{Step 3: Schwarz inequality.}
\begin{align}
    \int_t^{t+T} \norm{e(\tau)} \norm{\phi(\tau)} \, d\tau &\leq \left( \int_t^{t+T} \norm{e(\tau)}^2 \, d\tau \right)^{1/2} \left( \int_t^{t+T} \norm{\phi(\tau)}^2 \, d\tau \right)^{1/2}.
\end{align}

By PE, $\norm{\int_t^{t+T} \phi(\tau)^\top \phi(\tau) \, d\tau} \geq m\alpha$ for some $m > 0$, so $\int_t^{t+T} \norm{\phi(\tau)}^2 \, d\tau \geq c$ for some $c > 0$.

\textbf{Step 4: Exponential decay.}
The full proof requires showing that the PE condition forces $V(t)$ to decay exponentially. This is accomplished using Lyapunov-Krasovskii functional techniques (see \cite{Ioannou2006, Narendra1989}).

The key insight: PE prevents $\phi$ from becoming "small" over any interval, which combined with $e \in \Lp{2}$ and the feedback structure, forces both $e$ and $\thetatilde$ to zero.
\end{proof}

\subsection{Persistence of Excitation: Deeper Results}

\begin{theorem}[Spectral Characterization of PE]
\label{thm:pe_spectral}
Let $\phi(t) = \sum_{i=1}^n a_i \sin(\omega_i t + \psi_i)$ where $\omega_i$ are distinct frequencies. Then $\phi$ is PE if and only if $n \geq m$ (number of frequencies $\geq$ dimension of parameter space) and $a_i \neq 0$ for all $i$.
\end{theorem}

\begin{proof}[Proof Sketch]
The PE integral $\int_t^{t+T} \phi(\tau)^\top \phi(\tau) \, d\tau$ involves cross-terms $\sin(\omega_i t) \sin(\omega_j t)$. Over a period $T$:
\begin{itemize}
    \item If $\omega_i = \omega_j$, the integral is $\sim T/2$ (contributes to diagonal).
    \item If $\omega_i \neq \omega_j$, the integral averages to zero (orthogonality).
\end{itemize}

Thus, the matrix becomes approximately diagonal with entries proportional to $a_i^2 T/2$. For positive definiteness, we need all $a_i \neq 0$ and enough frequencies to fill the spectrum.

A rigorous proof involves Fourier analysis and the Riemann-Lebesgue lemma (see \cite{Sastry1989}).
\end{proof}

\section{Summary}

This chapter established:
\begin{itemize}
    \item Function spaces ($\Lp{p}$, Sobolev $W^{k,p}$) with examples relevant to adaptive control.
    \item Barbalat's Lemma with complete proof, plus the key corollary for $\Lp{2}$ signals.
    \item Persistent excitation definition, examples, and the fundamental theorem linking PE to parameter convergence.
\end{itemize}

These tools form the foundation for the Lyapunov stability analysis in subsequent chapters.
